% ----------------------------------------------------------------------------
%  3.3  Thực nghiệm trên dữ liệu chuẩn
% ----------------------------------------------------------------------------
\subsection{Tái hiện thí nghiệm trên dữ liệu Cosexp}\label{sec:cosexp}

Bộ dữ liệu Cosexp của hộp công cụ SVM-KM~\cite{canu2005} được bài báo gốc dùng để minh họa ảnh hưởng của số đặc trưng $M$. Các điểm được lấy đều trên miền $[0,5]\times[-1,1]$ và được gán nhãn theo dấu của $\cos\bigl(0{,}5\,e^{0{,}7x(1)}\bigr)-x(2)$; biên thật là đường cong $x(2)=\cos\bigl(0{,}5\,e^{0{,}7x(1)}\bigr)$, dao động ngày càng nhanh khi $x(1)$ tăng (nét đứt ở \cref{fig:cosexp}a,b). Theo đúng thiết lập của~\cite{huang2013}, ta dùng $500$ điểm huấn luyện, $500$ điểm kiểm tra và PWL-C-SVM với ánh xạ cộng tính, cho $M$ chạy từ $2$ đến $40$. Khác với bài báo, mỗi cấu hình được lặp lại trên $10$ bộ dữ liệu sinh độc lập, và ánh xạ HH được chạy song song để so sánh.

\begin{figure}[t]
\centering
\includegraphics[width=\linewidth]{fig_cosexp_M.pdf}
\caption[Tái hiện thí nghiệm Cosexp: ảnh hưởng của số đặc trưng]{Tái hiện thí nghiệm Cosexp (Hình~3 của~\cite{huang2013}), $10$ lần lặp. (a), (b) Biên của PWL-C-SVM cộng tính với $M=12$ và $M=40$ (nét liền) so với biên thật (nét đứt). (c) Độ chính xác trên tập kiểm tra theo $M$: đường là trung bình, dải màu là $\pm$ một độ lệch chuẩn; đường chấm là SVM hạt nhân Gauss. (d) Thời gian huấn luyện (trung vị) với $\gamma$ đã chọn.}
\label{fig:cosexp}
\end{figure}

\Cref{fig:cosexp} tái hiện đúng xu hướng mà bài báo gốc báo cáo. Với $M=2$, bộ phân loại chỉ là tuyến tính và đạt $72{,}7\pm2{,}0\%$ (bài báo: $76{,}83\%$ với một lần chạy). Khi $M$ tăng, biên uốn theo đường cong thật ngày càng sát (\cref{fig:cosexp}a,b); độ chính xác đạt $90{,}0\%$ tại $M=20=10n$ và ổn định quanh $94\%$ từ $M=24$, ngang với SVM hạt nhân Gauss ($94{,}2\pm0{,}9\%$). Ở $M$ nhỏ, đường cong không đơn điệu: mỗi giá trị của $M$ ứng với một lưới nút hoàn toàn khác, và độ chính xác phụ thuộc vào việc các nút có rơi gần những chỗ biên đổi hướng hay không. Thời gian huấn luyện tăng theo $M$ nhưng vẫn dưới $0{,}1$ giây (\cref{fig:cosexp}d).

Đáng chú ý, trên dữ liệu này ánh xạ HH kém hẳn ánh xạ cộng tính ($89{,}1\pm1{,}4\%$ tại $M=40$), dù về lý thuyết nó linh hoạt hơn. Lý do là nhãn của Cosexp được xác định bởi hàm $\cos\bigl(0{,}5\,e^{0{,}7x(1)}\bigr)-x(2)$, một hàm cộng tính: ánh xạ cộng tính khớp đúng cấu trúc đó, còn các siêu phẳng ngẫu nhiên của HH phân tán khả năng biểu diễn vào những hướng không cần thiết. Kết hợp với Mục~\ref{sec:exp2d}, ta rút ra một quy tắc thực hành: nên thử ánh xạ cộng tính trước, vì nó rẻ, dễ diễn giải và đủ tốt khi tương tác giữa các biến yếu; chỉ chuyển sang HH khi có dấu hiệu tương tác mạnh, như trên bàn cờ. Đây cũng là thứ tự mà bài báo gốc khuyến nghị~\cite[Mục~4]{huang2013}.

% Generated by code/fill_numbers.py from chuong3_bench_b.tpl -- edit the template, not this file.
% ----------------------------------------------------------------------------
\subsection{Mười bộ dữ liệu chuẩn}\label{sec:bench10}

\paragraph{Dữ liệu và giao thức.} Đề án dùng mười bộ dữ liệu phân loại nhị phân (\cref{tab:datasets}). Bảy bộ trong số đó có mặt trong Bảng~1 của bài báo gốc (Haberman, Transfusion, Pima, Breast, Magic, Parkinsons, Ionosphere); ba bộ còn lại (Banknote, Spambase, Sonar) được thêm vào để bao quát thêm các kích thước dữ liệu và số chiều khác nhau. Mỗi bộ được chia ngẫu nhiên có phân tầng thành hai nửa huấn luyện và kiểm tra, lặp lại $10$ lần; với Magic, tập huấn luyện gồm $2.000$ điểm như trong bài báo gốc, phần còn lại dùng để kiểm tra; với Spambase, do chi phí tinh chỉnh các mô hình hạt nhân lớn, chỉ lặp lại $5$ lần. Trên mỗi tập huấn luyện, tham số được chọn bằng kiểm định chéo $10$ phần: $\gamma\in\{2^{-5},2^{-3},\dots,2^{11}\}$ cho mọi mô hình SVM, tham số $1/\sigma^2$ của hạt nhân Gauss trên lưới $\{2^{-15},2^{-13},\dots,2^{3}\}$ theo khuyến nghị của~\cite{hsu2003}, và hệ số phạt $L_2$ của mạng nơ-ron trong $\{10^{-4},10^{-3},\dots,1\}$. Chín phương pháp được so sánh: SVM tuyến tính; SVM và LS-SVM hạt nhân Gauss; kNN với $k=1$ như trong bài báo; Ik-SVM; mạng nơ-ron một lớp ẩn gồm $D=9n$ nơ-ron ReLU, bằng đúng số đặc trưng bản lề của PWL-SVM dạng HH, được huấn luyện bằng L-BFGS~\cite{liu1989}; và ba biến thể PWL-SVM với $M/n=10$ theo đúng quy tắc của bài báo: PWL-C-SVM và PWL-LS-SVM với ánh xạ cộng tính lưới đều, và PWL-C-SVM với ánh xạ HH chuẩn hóa $|p(1)|=1$.

\begin{table}[t]
\centering
\caption[Mười bộ dữ liệu chuẩn]{Mười bộ dữ liệu chuẩn, sắp theo số chiều $n$ (sau khi bỏ các cột hằng). $N$ là số mẫu, cột cuối là số mẫu huấn luyện/kiểm tra trong mỗi lần chia.}
\label{tab:datasets}
\small
\renewcommand{\arraystretch}{1.15}
\setlength{\tabcolsep}{4.5pt}
\begin{tabular}{@{}l l r r r r@{}}
\toprule
\textbf{Tên} & \textbf{Nội dung} & $N$ & $n$ & \begin{tabular}[b]{@{}c@{}}\textbf{\% lớp}\\ \textbf{$+1$}\end{tabular} & \begin{tabular}[b]{@{}c@{}}\textbf{Huấn luyện/}\\ \textbf{kiểm tra}\end{tabular}\\
\midrule
Haberman & sống sót sau phẫu thuật ung thư vú & 306 & 3 & 26{,}5 & 153/153\\
Transfusion & hiến máu tình nguyện & 748 & 4 & 23{,}8 & 374/374\\
Banknote & xác thực tiền giấy & 1.372 & 4 & 44{,}5 & 686/686\\
Pima & bệnh tiểu đường & 768 & 8 & 34{,}9 & 384/384\\
Breast & ung thư vú Wisconsin (bản gốc) & 699 & 9 & 34{,}5 & 349/350\\
Magic & kính thiên văn tia gamma & 19.020 & 10 & 35{,}2 & 2.000/17.020\\
Parkinsons & bệnh Parkinson (giọng nói) & 195 & 22 & 75{,}4 & 97/98\\
Ionosphere & tín hiệu radar tầng điện li & 351 & 33 & 64{,}1 & 175/176\\
Spambase & thư rác & 4.601 & 57 & 39{,}4 & 2.300/2.301\\
Sonar & tín hiệu sonar (đá hay mìn) & 208 & 60 & 53{,}4 & 104/104\\
\bottomrule
\end{tabular}
\end{table}


\begin{table}[t]
\centering
\caption[Độ chính xác trên mười bộ dữ liệu chuẩn]{Độ chính xác trung bình trên tập kiểm tra (\%) qua $10$ lần chia ngẫu nhiên ($5$ lần với Spambase). In đậm: giá trị cao nhất mỗi hàng. Dòng cuối là hạng trung bình trên mười bộ dữ liệu ($1$ là tốt nhất). Độ lệch chuẩn được cho ở \cref{tab:appendix-bench} (Phụ lục~B).}
\label{tab:bench}
\footnotesize
\setlength{\tabcolsep}{3pt}
\renewcommand{\arraystretch}{1.18}
\resizebox{\linewidth}{!}{%
\begin{tabular}{@{}l r r r r r r r r r@{}}
\toprule
\textbf{Dữ liệu} & \begin{tabular}[b]{@{}c@{}}SVM\\tuyến tính\end{tabular} & \begin{tabular}[b]{@{}c@{}}SVM-\\RBF\end{tabular} & \begin{tabular}[b]{@{}c@{}}LS-SVM-\\RBF\end{tabular} & \begin{tabular}[b]{@{}c@{}}kNN\\$(k=1)$\end{tabular} & \begin{tabular}[b]{@{}c@{}}Ik-\\SVM\end{tabular} & \begin{tabular}[b]{@{}c@{}}MLP-\\ReLU\end{tabular} & \begin{tabular}[b]{@{}c@{}}PWL-C\\cộng tính\end{tabular} & \begin{tabular}[b]{@{}c@{}}PWL-LS\\cộng tính\end{tabular} & \begin{tabular}[b]{@{}c@{}}PWL-C\\HH\end{tabular}\\
\midrule
Haberman & 72{,}5 & 73{,}3 & 73{,}4 & 65{,}7 & 72{,}0 & 72{,}4 & 73{,}2 & 72{,}6 & \textbf{73{,}5}\\
Transfusion & 76{,}5 & 77{,}6 & 77{,}7 & 68{,}4 & 76{,}7 & \textbf{78{,}6} & 76{,}3 & 77{,}4 & 76{,}3\\
Banknote & 98{,}6 & \textbf{100{,}0} & \textbf{100{,}0} & 99{,}8 & 99{,}9 & 99{,}9 & 99{,}9 & 99{,}5 & 99{,}9\\
Pima & \textbf{76{,}7} & 76{,}1 & 76{,}6 & 70{,}8 & 76{,}0 & \textbf{76{,}7} & 76{,}3 & 76{,}5 & 76{,}2\\
Breast & 96{,}6 & 96{,}5 & 96{,}4 & 95{,}3 & 96{,}4 & \textbf{96{,}8} & 96{,}3 & 96{,}0 & 95{,}6\\
Magic & 79{,}0 & 85{,}5 & 85{,}6 & 78{,}3 & 85{,}0 & \textbf{85{,}7} & 84{,}8 & 84{,}4 & 84{,}3\\
Parkinsons & 85{,}4 & 89{,}9 & 91{,}4 & \textbf{93{,}7} & 89{,}4 & 90{,}6 & 86{,}2 & 89{,}3 & 85{,}5\\
Ionosphere & 86{,}2 & 93{,}4 & \textbf{93{,}8} & 86{,}4 & 91{,}1 & 88{,}0 & 89{,}7 & 87{,}7 & 87{,}0\\
Spambase & 93{,}0 & 93{,}6 & 93{,}4 & 88{,}8 & \textbf{94{,}6} & 93{,}9 & 93{,}4 & 92{,}6 & 82{,}1\\
Sonar & 75{,}2 & 84{,}3 & \textbf{84{,}5} & 81{,}2 & 82{,}0 & 81{,}9 & 80{,}2 & 79{,}8 & 76{,}2\\
\midrule
Hạng trung bình & 6{,}55 & 3{,}15 & 2{,}35 & 7{,}40 & 4{,}60 & 3{,}00 & 5{,}35 & 6{,}00 & 6{,}60\\
\bottomrule
\end{tabular}}
\end{table}


\paragraph{Kết quả tổng quát.} \Cref{tab:bench} cho độ chính xác trung bình và \cref{fig:bench-diff} biểu diễn chênh lệch so với SVM hạt nhân Gauss; độ lệch chuẩn và kích thước mô hình được cho ở Phụ lục~B. Ba phương pháp dẫn đầu về độ chính xác là LS-SVM hạt nhân Gauss, mạng nơ-ron và SVM hạt nhân Gauss, với hạng trung bình lần lượt $2{,}35$, $3{,}00$ và $3{,}15$. PWL-C-SVM cộng tính đứng giữa bảng (hạng $5{,}35$), trên SVM tuyến tính ($6{,}55$) và kNN ($7{,}40$).

Bức tranh rõ hơn khi tách các bộ dữ liệu thành hai nhóm. Ở nhóm thứ nhất gồm Haberman, Transfusion, Banknote, Pima và Breast, mọi phương pháp dựa trên SVM, kể cả SVM tuyến tính, chênh nhau không quá khoảng $1{,}5$ điểm phần trăm: một biên tuyến tính đã gần như đủ tốt, và tính phi tuyến không đem lại lợi ích đáng kể. Ở nhóm thứ hai gồm Magic, Parkinsons, Ionosphere, Sonar, nơi SVM tuyến tính kém SVM hạt nhân Gauss từ $4$ đến $9$ điểm, PWL-C-SVM cộng tính thu hẹp đáng kể khoảng cách trên ba trong bốn bộ: nó cải thiện so với SVM tuyến tính $5{,}8$ điểm trên Magic, $3{,}5$ điểm trên Ionosphere và $5{,}0$ điểm trên Sonar, nhưng vẫn kém SVM hạt nhân Gauss từ $0{,}7$ đến $4{,}1$ điểm. Magic là trường hợp đáng chú ý nhất: PWL-C-SVM đạt $84{,}8\%$ so với $85{,}5\%$ của SVM hạt nhân Gauss, trong khi mô hình chỉ gồm $211$ số thực so với $7.694$ số, nhỏ hơn khoảng $36$ lần.

\begin{figure}[t]
\centering
\includegraphics[width=\linewidth]{fig_benchmark_diff.pdf}
\caption[Chênh lệch độ chính xác so với SVM hạt nhân Gauss]{Chênh lệch độ chính xác trung bình so với SVM hạt nhân Gauss (điểm phần trăm; giá trị âm nghĩa là kém hơn) của bốn phương pháp trên mười bộ dữ liệu.}
\label{fig:bench-diff}
\end{figure}

\begin{table}[t]
\centering
\caption[So sánh theo cặp với PWL-C-SVM cộng tính]{So sánh theo cặp giữa PWL-C-SVM cộng tính và từng phương pháp khác trên mười bộ dữ liệu: số bộ mà PWL-C-SVM cao hơn / bằng / thấp hơn (theo độ chính xác trung bình làm tròn đến $0{,}1$ điểm phần trăm), chênh lệch trung bình, và giá trị $p$ của kiểm định dấu hạng Wilcoxon hai phía.}
\label{tab:wilcoxon}
\small
\renewcommand{\arraystretch}{1.18}
\begin{tabular}{@{}l c r r@{}}
\toprule
\textbf{Phương pháp so sánh} & \textbf{Cao hơn / bằng / thấp hơn} & \textbf{Chênh lệch TB} & $p$\\
\midrule
SVM tuyến tính & 7 / 0 / 3 & $+1{,}66$ & 0{,}037\\
SVM-RBF & 1 / 0 / 9 & $-1{,}39$ & 0{,}014\\
LS-SVM-RBF & 0 / 1 / 9 & $-1{,}65$ & 0{,}008\\
kNN ($k=1$) & 8 / 0 / 2 & $+2{,}79$ & 0{,}105\\
Ik-SVM & 2 / 1 / 7 & $-0{,}68$ & 0{,}110\\
MLP-ReLU & 2 / 1 / 7 & $-0{,}82$ & 0{,}154\\
PWL-LS-SVM cộng tính & 7 / 0 / 3 & $+0{,}05$ & 0{,}432\\
PWL-C-SVM HH & 7 / 2 / 1 & $+1{,}97$ & 0{,}025\\
\bottomrule
\end{tabular}
\end{table}


\paragraph{Kiểm định thống kê.} \Cref{tab:wilcoxon} so sánh PWL-C-SVM cộng tính với từng phương pháp bằng kiểm định dấu hạng Wilcoxon trên mười cặp giá trị trung bình, theo khuyến nghị của Demšar~\cite{demsar2006}. Ở mức ý nghĩa $5\%$, PWL-C-SVM cộng tính kém SVM hạt nhân Gauss và LS-SVM hạt nhân Gauss một cách có ý nghĩa ($p=0{,}014$ và $p=0{,}008$). Nó cao hơn SVM tuyến tính trên bảy trong mười bộ, với chênh lệch trung bình $1{,}7$ điểm, và khác biệt này có ý nghĩa thống kê ($p=0{,}037$); khác biệt so với mạng nơ-ron và Ik-SVM không có ý nghĩa thống kê. Với chỉ mười bộ dữ liệu, các kiểm định này có lực thấp, nên “không bác bỏ” không có nghĩa là “tương đương”.

\paragraph{Ánh xạ HH và quy tắc chọn tham số.} PWL-C-SVM với ánh xạ HH có hạng trung bình thấp ($6{,}60$), và trên Spambase kết quả của nó rất bất ổn: ở lần chia đầu tiên, mô hình thất bại hoàn toàn ($40{,}0\%$, thấp hơn cả tỉ lệ $60{,}6\%$ của lớp đa số), trong khi ở bốn lần chia còn lại nó đạt từ $91{,}0\%$ đến $93{,}7\%$. Nguyên nhân nằm ở quy tắc sinh tham số chứ không ở mô hình. Dữ liệu Spambase rất thưa: sau khi chuẩn hóa, gần $78\%$ các giá trị bằng $0$, và phần lớn dữ liệu nằm sát một đỉnh của hình hộp $[0,1]^{57}$. Các siêu phẳng đi qua những điểm lấy đều trong hình hộp vì vậy thường không cắt qua vùng dữ liệu: ở mỗi lần chia, từ 18\% đến 22\% số đặc trưng bản lề bằng $0$ trên mọi điểm huấn luyện, và từ 17\% đến 22\% luôn hoạt động, tức là chỉ tuyến tính trên dữ liệu. Nghiêm trọng hơn, việc cố định $|p(1)|=1$ (\cref{rem:scale}) làm độ lớn của các đặc trưng phụ thuộc mạnh vào lần sinh ngẫu nhiên: giá trị lớn nhất của các đặc trưng trên tập huấn luyện dao động từ $423$ đến $3{,}3\cdot10^{5}$ giữa các lần chia. Ở lần chia đầu tiên, một siêu phẳng gần song song với trục $x(1)$ có hệ số lên tới $5{,}3\cdot10^{4}$, khiến bài toán tối ưu bị điều kiện hóa rất kém và bộ giải LIBLINEAR không hội tụ tới một nghiệm có ý nghĩa. Chỉ cần chuẩn hóa $\|p\|_2=1$, các đặc trưng trên dữ liệu huấn luyện không vượt quá $2{,}0$ ở mọi lần chia, và với cùng giao thức, PWL-C-SVM HH đạt $93{,}4\pm0{,}2\%$ trên Spambase. Nhìn chung, mạng nơ-ron với cùng số nơ-ron ẩn tốt hơn PWL-C-SVM HH trên tám trong mười bộ: học lớp ẩn từ dữ liệu có lợi hơn nhiều so với sinh nó ngẫu nhiên, và đó là cái giá của tính lồi đã được nêu ở Mục~\ref{sec:pwlsvm}.

\paragraph{So với bài báo gốc.} Kết quả trên không xác nhận nhận định của bài báo gốc rằng PWL-SVM có độ chính xác tương đương SVM hạt nhân Gauss~\cite{huang2013}. Nguyên nhân chính là bài báo chỉ dùng một lần chia dữ liệu: trên các bộ nhỏ như Parkinsons, với $98$ điểm kiểm tra, độ lệch chuẩn giữa các lần chia lên tới $2$--$4$ điểm phần trăm, nên một lần chia có thể cho kết quả rất lạc quan; bài báo báo cáo độ chính xác $100\%$ cho PWL-C-SVM trên bộ này, trong khi trung bình qua mười lần chia là $86{,}2\%$. Mặt khác, lợi thế về kích thước mô hình được xác nhận rõ ràng: với ánh xạ cộng tính, PWL-C-SVM cần từ $64$ đến $1.261$ số thực, nhỏ hơn SVM hạt nhân Gauss từ $2{,}5$ đến $36$ lần trên cả mười bộ dữ liệu (Phụ lục~B).

% ----------------------------------------------------------------------------
\subsection{Ảnh hưởng của số đặc trưng và của quy tắc chọn tham số}\label{sec:sensitivity}

Để trả lời câu hỏi CH3 một cách có kiểm soát, ta cố định mô hình là PWL-LS-SVM, giải chính xác bằng hồi quy ridge (\cref{prop:ridge}) để loại trừ ảnh hưởng của bộ giải, và thay đổi hai yếu tố: số đặc trưng trên mỗi chiều $M/n\in\{1,2,3,5,10,20\}$ và quy tắc chọn tham số. Bốn quy tắc được so sánh: ánh xạ cộng tính với lưới đều (quy tắc của bài báo gốc) và với lưới phân vị (\cref{rem:quantile}); ánh xạ HH với $|p(1)|=1$ (quy tắc của bài báo gốc) và với $\|p\|_2=1$ (\cref{rem:scale}). Thí nghiệm được thực hiện trên ba bộ dữ liệu có đặc điểm khác nhau --- Ionosphere ($n=33$), Magic ($n=10$, nhiều dữ liệu) và Spambase ($n=57$, rất thưa) --- với $5$ lần chia như ở Mục~\ref{sec:bench10} và $\gamma$ chọn bằng kiểm định chéo $5$ phần. Kết quả được vẽ ở \cref{fig:sensitivity}.

\begin{figure}[t]
\centering
\includegraphics[width=\linewidth]{fig_sensitivity.pdf}
\caption[Ảnh hưởng của số đặc trưng và quy tắc chọn tham số]{Độ chính xác trung bình (qua $5$ lần chia) của PWL-LS-SVM theo số đặc trưng trên mỗi chiều $M/n$, với bốn quy tắc chọn tham số. $M/n=1$ ứng với bộ phân loại tuyến tính.}
\label{fig:sensitivity}
\end{figure}

Có bốn nhận xét. Thứ nhất, độ chính xác tăng nhanh khi $M/n$ tăng từ $1$ (bộ phân loại tuyến tính) và bão hòa quanh $M/n=5$--$10$, phù hợp với quy tắc $M=10n$ của bài báo gốc; tăng tiếp lên $M/n=20$ hầu như không có lợi. Thứ hai, đặt nút theo phân vị cho kết quả tốt hơn hoặc ngang lưới đều trên cả ba bộ dữ liệu, và khác biệt rõ nhất trên những dữ liệu có phân bố lệch: trên Magic, lưới phân vị đạt $84{,}2\%$ ngay với $M/n=2$, trong khi lưới đều cần $M/n=10$ để đạt $84{,}3\%$; trên Spambase, lưới phân vị đạt $93{,}9\%$ so với $92{,}6\%$ của lưới đều ở $M/n=10$, đồng thời chỉ dùng $156$ đặc trưng thay vì $570$, vì các nút trùng nhau của những biến gần như luôn bằng $0$ bị loại bỏ. Thứ ba, ánh xạ HH cần nhiều đặc trưng hơn hẳn mới đạt tới độ chính xác của ánh xạ cộng tính: trên Magic, nó chỉ đuổi kịp lưới phân vị khi $M/n\ge10$ (và nhỉnh hơn $0{,}4$ điểm ở $M/n=20$), còn trên Spambase nó kém ở mọi giá trị của $M/n$; chỉ trên Ionosphere, nơi tương tác giữa các biến có vai trò, HH với $\|p\|_2=1$ đạt kết quả tốt nhất ($90{,}0\%$ ở $M/n=10$). Thứ tư, chuẩn hóa $\|p\|_2=1$ không bao giờ làm kết quả xấu đi đáng kể và cải thiện rõ trên Ionosphere ($90{,}0\%$ so với $86{,}9\%$ ở $M/n=10$); kết hợp với sự cố của LIBLINEAR trên Spambase, đây là một lý do mạnh để thay quy tắc $|p(1)|=1$ bằng $\|p\|_2=1$.

Từ đó rút ra một quy trình thực hành cho dữ liệu dạng bảng: bắt đầu với ánh xạ cộng tính, nút theo phân vị và $M/n$ khoảng $5$ đến $10$; chỉ dùng ánh xạ HH, với các siêu phẳng được chuẩn hóa, khi có dấu hiệu tương tác mạnh giữa các biến. Quy trình này được áp dụng cho bài toán tín dụng ở Mục~\ref{sec:application}.

